\chapter{Archiviazione e governance con MongoDB Time Series}

\section{Layer Gold e archiviazione con MongoDB Time Series}
Il livello \textbf{Gold} rappresenta lo strato di vertice dell'architettura dati, destinato ad alimentare in modo tempestivo e strutturato le interrogazioni analitiche, l'interfaccia grafica utente e gli algoritmi di rilevamento delle anomalie. I dati bonificati nel livello \textbf{Silver} vengono convogliati verso una base di dati MongoDB versione 8.0, all'interno della quale è configurata una collezione temporale nativa denominata \texttt{vital\_signals}.

La funzionalità \textbf{Time Series Collection}, introdotta nelle recenti versioni di MongoDB, costituisce un'innovazione architetturale di rilievo rispetto alle collezioni a documenti tradizionali. In un database a documenti convenzionale, la memorizzazione di milioni di misurazioni ad alta frequenza comporterebbe un overhead insostenibile: ogni singolo secondo richiederebbe un documento BSON completo con il rispettivo identificatore univoco \texttt{\_id}, le chiavi di campo replicate e le strutture di indice correlate. Nelle collezioni temporali, invece, MongoDB astrae la gestione fisica del dato raggruppando in modo del tutto trasparente le osservazioni temporali contigue all'interno di blocchi compressi denominati \textbf{bucket}.

La creazione della collezione \texttt{vital\_signals}, implementata nello script \texttt{load\_mongo.py}, richiede la definizione esplicita di tre parametri architetturali:
il parametro \texttt{timeField}, impostato sul campo radice \texttt{timestamp}, che indica al motore il campo di riferimento temporale per ordinare e partizionare i record;
il parametro \texttt{metaField}, impostato sul sotto-documento \texttt{metadata}, che raccoglie gli attributi descrittivi immutabili della serie associata a quel paziente;
il parametro \texttt{granularity}, impostato al valore \texttt{seconds}, che suggerisce al motore la frequenza attesa delle misurazioni per ottimizzare la dimensione e la durata temporale di ciascun \textbf{bucket}.

All'interno dell'infrastruttura di archiviazione, MongoDB organizza i dati nella collezione di sistema interna \texttt{system.buckets.vital\_signals}. Quando giungono nuovi punti recanti i medesimi metadati, il motore inserisce i campioni nel \textbf{bucket} aperto corrente fino al raggiungimento della capacità massima di compressione (fino a 1.000 punti o un limite temporale stabilito). Per comprimere i dati all'interno del \textbf{bucket}, il motore applica tecniche colonnari avanzate: i timestamp crescenti vengono codificati mediante l'algoritmo \textbf{Delta-of-Delta}, che memorizza unicamente le variazioni degli incrementi temporali riducendo l'ingombro a pochi bit per punto, mentre i vettori delle metriche numeriche vengono compressi tramite l'algoritmo \textbf{Zstandard}. In questo modo viene massimizzata la densità informativa ed eliminata qualsiasi ridondanza dei metadati contestuali.

\section{Benchmark prestazionale e di occupazione dello storage}
Per valutare in modo quantitativo e scientificamente rigoroso i benefici apportati dalla transizione dai file Parquet alla collezione temporale di MongoDB, è stata eseguita una campagna di benchmark automatizzata mediante lo script \texttt{benchmark\_etl.py}.

La misurazione ha preso in esame un campione a larga scala di 50 casi clinici completi, per un totale di 311.173 punti temporali bonificati. La sperimentazione ha indagato tre parametri cardine: l'occupazione fisica dello storage su disco, la latenza di scansione delle interrogazioni temporali e la completezza delle celle numeriche. I risultati ottenuti sono compendiati nella Tabella~\ref{tab:etl_benchmark} e illustrati graficamente nella Figura~\ref{fig:etl_benchmark}.

\begin{table}[ht]
\centering
\resizebox{\textwidth}{!}{%
\begin{tabular}{lccc}
\toprule
\textbf{Parametro / Metrica} & \textbf{Bronze (grezzo)} & \textbf{Silver (pulito Parquet)} & \textbf{Gold (MongoDB TS)} \\
\midrule
Occupazione disco (50 casi) & 3.88 MB & 2.24 MB (-42.4\%) & 0.42 MB (-89.1\%) \\
Latenza scansione singolo caso & 13.48 ms (Parquet) & 13.48 ms (Parquet) & 3.63 ms (Pushdown indicizzato) \\
Fattore di velocizzazione & 1.0$\times$ (Riferimento) & 1.0$\times$ & \textbf{3.7$\times$ più veloce} \\
Valori mancanti misurati & 67.6\% & 32.6\% (post-LOCF) & 0.0\% (Omissione BSON) \\
\bottomrule
\end{tabular}}
\caption{Confronto quantitativo di storage, latenza e qualità dati tra i layer su 311.173 record.}
\label{tab:etl_benchmark}
\end{table}

\begin{figure}[ht]
    \centering
    % Diagramma TikZ comparativo per Storage e Latenza di Query (Benchmark ETL)
\begin{tikzpicture}[
    scale=0.92, transform shape,
    font=\sffamily,
    cardBox/.style={
        draw=brandBorder,
        fill=white,
        rounded corners=8pt,
        line width=1pt,
        inner sep=0pt,
        drop shadow={opacity=0.12, shadow xshift=2pt, shadow yshift=-2pt}
    },
    bgBar/.style={
        fill=gray!12,
        rounded corners=3pt
    },
    valBar/.style={
        rounded corners=3pt
    }
]

% ==================== CARD 1: STORAGE SU DISCO ====================
\node[cardBox, minimum width=7.2cm, minimum height=6.8cm] (storageCard) at (0, 0) {};

% Titolo Card 1
\node[font=\bfseries\small, text=brandDark, anchor=north] at (0, 3.1) {Occupazione storage su disco};
\node[font=\scriptsize, text=brandMuted, anchor=north] at (0, 2.65) {Dimensione per 500 casi clinici (5.402.386 record)};

% --- Riga 1: Bronze ---
\node[font=\scriptsize, text=brandCharcoal, anchor=west] at (-3.2, 1.95) {\textbf{Layer Bronze} \textcolor{brandMuted}{(Grezzo)}};
\node[font=\scriptsize, text=brandCoral!90!black, anchor=east] at (3.2, 1.95) {\textbf{33.95 MB} \textcolor{brandMuted}{(6 col.)}};
\draw[bgBar] (-3.2, 1.45) rectangle (3.2, 1.75);
\draw[fill=brandCoral, valBar] (-3.2, 1.45) rectangle (1.68, 1.75);

% --- Riga 2: Silver ---
\node[font=\scriptsize, text=brandCharcoal, anchor=west] at (-3.2, 0.95) {\textbf{Layer Silver} \textcolor{brandMuted}{(18 col. DQ)}};
\node[font=\scriptsize, text=brandPrimaryDark, anchor=east] at (3.2, 0.95) {\textbf{44.52 MB} \textcolor{brandMuted}{(+31.1\%)}};
\draw[bgBar] (-3.2, 0.45) rectangle (3.2, 0.75);
\draw[fill=brandPrimary!80, valBar] (-3.2, 0.45) rectangle (3.2, 0.75);

% --- Riga 3: Gold ---
\node[font=\scriptsize, text=brandCharcoal, anchor=west] at (-3.2, -0.05) {\textbf{Layer Gold} \textcolor{brandMuted}{(Mongo TS puro)}};
\node[font=\scriptsize, text=brandTeal!90!black, anchor=east] at (3.2, -0.05) {\textbf{15.60 MB} \textcolor{brandTeal}{\textbf{(-54.1\% / -65\% vs S)}}};
\draw[bgBar] (-3.2, -0.55) rectangle (3.2, -0.25);
\draw[fill=brandTeal, valBar] (-3.2, -0.55) rectangle (-0.96, -0.25);

% Callout riassuntivo in basso
\node[fill=brandTeal!12, draw=brandTeal!50, line width=0.8pt, rounded corners=5pt, inner sep=6pt, font=\scriptsize\bfseries, text=brandTeal!90!black] at (0, -2.1) {
    Storage puro: 15.60 MB (+ 2.56 MB indici B-Tree) $\cdot$ -65.0\% vs Silver
};


% ==================== CARD 2: LATENZA QUERY ====================
\node[cardBox, minimum width=7.2cm, minimum height=6.8cm] (latencyCard) at (8.0, 0) {};

% Titolo Card 2
\node[font=\bfseries\small, text=brandDark, anchor=north] at (8.0, 3.1) {Latenza scansione query};
\node[font=\scriptsize, text=brandMuted, anchor=north] at (8.0, 2.65) {Mediana su 15 esecuzioni con warm-up (caso 1)};

% --- Riga 1: Parquet ---
\node[font=\scriptsize, text=brandCharcoal, anchor=west] at (4.8, 1.85) {\textbf{Parquet partizionato}};
\node[font=\scriptsize, text=brandCoral!90!black, anchor=east] at (11.2, 1.85) {\textbf{26.95 ms} \textcolor{brandMuted}{(1.0$\times$)}};
\draw[bgBar] (4.8, 1.35) rectangle (11.2, 1.65);
\draw[fill=brandCoral, valBar] (4.8, 1.35) rectangle (11.2, 1.65);

% --- Riga 2: MongoDB pushdown ---
\node[font=\scriptsize, text=brandCharcoal, anchor=west] at (4.8, 0.75) {\textbf{MongoDB Gold}};
\node[font=\scriptsize, text=brandPrimaryDark, anchor=east] at (11.2, 0.75) {\textbf{12.02 ms} \textcolor{brandPrimary}{\textbf{(2.2$\times$ speedup)}}};
\draw[bgBar] (4.8, 0.25) rectangle (11.2, 0.55);
\draw[fill=brandPrimary, valBar] (4.8, 0.25) rectangle (7.65, 0.55);

% Callout riassuntivo in basso
\node[fill=brandPrimary!8, draw=brandPrimary!40, line width=0.8pt, rounded corners=5pt, inner sep=6pt, font=\scriptsize, align=center, text=brandDark] at (8.0, -1.95) {
    \textbf{\textcolor{brandPrimary}{2.2$\times$ Più Veloce}} \quad $\cdot$ \quad \textcolor{brandDark!80}{Pushdown indici su metadati di bucket}
};

\end{tikzpicture}

    \caption{Benchmark comparativo dell'occupazione disco e della latenza di query generato in TikZ.}
    \label{fig:etl_benchmark}
\end{figure}

Come emerge dai dati riportati nella Tabella~\ref{tab:etl_benchmark}, il passaggio dai file Parquet non compressi del livello \textbf{Bronze} (3.88 MB) al formato bonificato del livello \textbf{Silver} (2.24 MB) produce una prima contrazione del volume del 42.4\%, dovuta alla rimozione delle righe completamente prive di segnale vitale. L'ulteriore transizione al livello \textbf{Gold} su MongoDB abbatte lo spazio occupato a soli 0.42 MB, facendo registrare un risparmio complessivo di storage dell'89.1\% rispetto al dato grezzo originario. Questo risultato dimostra l'eccellente efficienza dell'algoritmo di compressione a \textbf{bucket} per serie temporali biometriche a elevata frequenza.

Sul fronte delle prestazioni computazionali, il protocollo di test ha adottato una metodologia comparativa uniforme: è stata misurata la mediana dei tempi di risposta su quindici iterazioni successive relative al caso clinico 1, precedute da una fase di riscaldamento (\textbf{warmup}) per stabilizzare le cache del sistema operativo e del database. L'interrogazione mirava a estrarre la media della frequenza cardiaca per l'intero arco temporale dell'intervento. La lettura del file Parquet partizionato ha richiesto una latenza mediana di 13.48 ms, mentre la corrispondente query di aggregazione in MongoDB, eseguita tramite pipeline con operatore \texttt{\$match} e \texttt{\$group}, è stata evasa in appena 3.63 ms. Il motore di MongoDB consegue dunque un fattore di velocizzazione pari a 3.7 volte, beneficiando del filtraggio anticipato a livello di metadati del \textbf{bucket} prima di procedere alla decodifica delle serie numeriche.

\section{Data governance artigianale in ottica data lakehouse}
Un aspetto centrale del progetto, concordato e approfondito durante le fasi di revisione con la relatrice, risiede nella formalizzazione della \textbf{data governance}. Nei moderni contesti industriali, l'adozione di un'architettura a livelli richiede meccanismi stringenti per garantire l'integrità strutturale, il tracciamento della provenienza e l'affidabilità delle transazioni. Piattaforme dedicate quali Delta Lake, Apache Iceberg o Apache Hudi forniscono queste garanzie mediante registri di transazioni distribuiti (\textbf{transaction log}) e cataloghi integrati all'interno di motori di calcolo analitico quali Spark o Trino.

Nel presente lavoro, per mantenere l'architettura compatta, autosufficiente ed eseguibile in ambienti ospedalieri locali senza dipendenze computazionali onerose, si è scelto di emulare artigianalmente le proprietà fondamentali di un \textbf{data lakehouse} sfruttando esclusivamente gli strumenti e le primitive native di MongoDB. Tale impostazione poggia su quattro componenti complementari: il catalogo di registrazione delle serie, la tabella dei metadati descrittivi, l'applicazione di schemi formali di validazione e una strategia applicativa di compensazione per le scritture.

\subsection{Catalogo e registrazione della provenienza (collezione registry)}
Per tenere traccia dello stato di ingestione di ciascun caso clinico e garantire la ricostruibilità storica del processo di caricamento, è stata istituita la collezione \texttt{registry}. Tale collezione funge da catalogo centrale e opera come l'equivalente di un metastore applicativo.

Ogni qualvolta una serie temporale viene promossa dal livello \textbf{Silver} al livello \textbf{Gold}, viene inserito nel registro un documento che attesta l'avvenuta pubblicazione. Come mostrato nel Listato~\ref{lst:registry_schema}, il documento registra l'identificativo univoco del caso clinico (\texttt{case\_id}), il numero complessivo di punti temporali validati e immessi nella serie (\texttt{record\_count}), la versione formale dello schema logico adottato (\texttt{schema\_version}) e l'oggetto di provenienza (\texttt{provenance}), che riporta la fase di elaborazione (\texttt{silver\_to\_gold}) e l'istante temporale esatto di caricamento con marcatura UTC.

\begin{lstlisting}[caption={Schema formale di validazione JSON Schema per la collezione registry.}, label={lst:registry_schema}]
// Validazione applicata alla collezione 'registry' in MongoDB
db.runCommand({
  collMod: "registry",
  validator: {
    $jsonSchema: {
      bsonType: "object",
      required: ["case_id", "schema_version", "provenance"],
      properties: {
        case_id: {
          bsonType: ["int", "long"],
          description: "ID del caso clinico, intero a 32 o 64 bit obbligatorio"
        },
        schema_version: {
          bsonType: "string",
          description: "Versione dello schema, stringa obbligatoria"
        },
        record_count: {
          bsonType: ["int", "long"],
          description: "Numero complessivo di record temporali caricati"
        },
        provenance: {
          bsonType: "object",
          required: ["step", "timestamp"],
          properties: {
            step: { bsonType: "string" },
            timestamp: { bsonType: "date" }
          }
        }
      }
    }
  },
  validationLevel: "strict",
  validationAction: "error"
});
\end{lstlisting}

La collezione \texttt{registry} è protetta da regole native di \textbf{JSON Schema Validation} configurate con livello rigoroso (\texttt{validationLevel: "strict"}) e azione di rifiuto immediato (\texttt{validationAction: "error"}). In questo modo, qualsiasi tentativo di inserire nel catalogo record incompleti o privi dei vincoli di provenienza viene categoricamente respinto dal motore del database.

\subsection{Metadati descrittivi delle tracce (collezione track\_metadata)}
Accanto al catalogo delle esecuzioni, la governance necessita di una definizione semantica chiara del significato di ciascuna variabile misurata. A questo scopo è stata creata la collezione ausiliaria \texttt{track\_metadata}, inizializzata tramite lo script \texttt{mongo/registry\_schema.js}.

Per ciascuna metrica fisiologica registrata dal monitor \texttt{Solar8000}, la collezione documenta il nome formale della traccia, la versione del contratto dati, l'unità di misura fisica adottata (\texttt{bpm}, percentuale o \texttt{mmHg}) e l'intervallo numerico fisiologicamente plausibile (\texttt{expected\_range}). Questa collezione funge da vero e proprio contratto di interfaccia tra il personale sanitario, i moduli di bonifica e i consumatori a valle, chiarendo in modo inequivocabile i limiti attesi di ogni segnale.

\subsection{Limiti transazionali e strategia applicativa di compensazione}
Un'analisi rigorosa della piattaforma impone di evidenziare con trasparenza i vincoli intrinseci dell'infrastruttura scelta. Come documentato nel manuale ufficiale di MongoDB, le \textbf{Time Series Collection} non consentono operazioni di inserimento all'interno di transazioni multi-documento distribuite; l'esecuzione di una scrittura temporale sotto una sessione transazionale genera l'errore nativo 263 (\texttt{OperationNotSupportedInTransaction}). Ne consegue che non è tecnicamente possibile effettuare un'operazione atomica ACID congiunta tra la scrittura fisica delle serie in \texttt{vital\_signals} e l'aggiornamento del catalogo in \texttt{registry}.

Per ovviare a tale limite senza compromettere la coerenza dei dati, nel modulo di caricamento \texttt{load\_mongo.py} è stata adottata una strategia applicativa di compensazione ordinata:
la procedura esegue dapprima l'inserimento in blocco di tutti i punti della serie all'interno della collezione \texttt{vital\_signals};
solo a fronte del completamento con esito positivo dell'inserimento dei dati temporali, la procedura procede alla registrazione del caso nel catalogo \texttt{registry};
qualora la prima operazione dovesse sollevare un errore di rete o di I/O, il controllo interrompe l'esecuzione e il documento di catalogo non viene scritto. In questo modo si impedisce l'insorgenza di casi fantasma all'interno del catalogo.
Al contempo, per rendere l'operazione idempotente e prevenire duplicazioni accidentali in caso di riavvio della procedura, il caricatore verifica preliminarmente la presenza del caso nel \texttt{registry}: se il caso risulta già catalogato, il caricamento dei dati associati viene saltato a priori.

\section{Uso dei metadati per la data quality e analisi critica}
Nel disegno dell'architettura, i metadati svolgono una duplice funzione: abilitano prestazioni di accesso ottimali e governano l'osservabilità della qualità del dato lungo tutte le fasi della pipeline.

Il campo \texttt{metaField} della collezione temporale svolge un ruolo tecnologico primario: raccogliendo in modo strutturato l'identificativo del caso, il monitor di provenienza, il reparto chirurgico, l'età e il sesso del paziente, MongoDB genera automaticamente indici composti che accoppiano i metadati al timestamp. Le interrogazioni analitiche che filtrano per paziente o per reparto beneficiano del cosiddetto \textbf{query pushdown}: il motore individua ed esamina esclusivamente i \textbf{bucket} pertinenti senza dover decodificare le serie storiche dei pazienti estranei alla query.

Inoltre, per garantire che le valutazioni di qualità effettuate nel livello \textbf{Silver} non vadano disperse nella successiva fase di analisi o addestramento dei modelli, ciascun punto memorizzato nel livello \textbf{Gold} incorpora l'oggetto \texttt{quality\_flags}. In tal modo, un algoritmo di \textbf{machine learning} o un clinico possono interrogare la serie distinguendo immediatamente le misurazioni fisiologicamente conformi da quelle che avevano evidenziato valori anomali durante i controlli di plausibilità.

A supporto della tracciabilità complessiva, la conclusione della fase di pulizia \textbf{Silver} comporta la generazione automatica del file \texttt{quality\_report.json}, che censisce per ogni singolo paziente il numero di campioni originali, i campioni mantenuti, le righe rimosse, il tasso di valori nulli residui e il conteggio analitico degli outlier rilevati per ciascun canale. Al fine di mantenere memoria storica dell'evoluzione della qualità, la pipeline provvede ad archiviare la copia precedente del report all'interno della cartella \texttt{quality\_history} assegnando una marca temporale univoca.

A fronte di questi risultati, è doveroso svolgere alcune considerazioni critiche sulla semantica delle metriche rilevate:
la completezza dichiarata pari a zero valori nulli nel livello \textbf{Gold} (come evidenziato nella Tabella~\ref{tab:etl_benchmark}) è la conseguenza diretta della struttura documentale BSON e dell'eliminazione preventiva delle righe totalmente vuote. In MongoDB, i campi che non presentano una misura valida vengono omessi dal sotto-documento \texttt{metrics} per non occupare spazio; tale assenza non deve essere confusa con un recupero matematico miracoloso di tutti i dati mancanti originali, ma rappresenta un'ottimizzazione rappresentativa del dato;
il processo di bonifica nel livello \textbf{Silver} elimina 333.464 righe su 644.637 campioni complessivi presenti nel \textbf{Bronze}. Questa operazione abbatte drasticamente il rumore e dimezza il tasso di celle nulle (dal 67.6\% al 32.6\%), ma comporta la perdita di oltre la metà dei timestamp iniziali, corrispondenti alle fasi in cui i sensori non registravano segnali utili;
l'asse temporale è stato normalizzato sommando i secondi relativi dell'intervento a una data convenzionale fissa (\texttt{2020-01-01 + Time}). Sebbene tale espediente risponda pienamente alle esigenze di ordinamento e indicizzazione delle collezioni temporali, in una futura implementazione ospedaliera di produzione sarà opportuno raccordare i campioni all'orario effettivo registrato nel cartellino anestesiologico.
